\chapter*{考研三卷种阅读路线}
\addcontentsline{toc}{chapter}{考研三卷种阅读路线}
本书覆盖面宽于统考所需内容。数学二不设置概率论与数理统计部分；数学一、三需要按各自范围阅读，不能将本书第1—10章都视为共同必修。大纲中的“考试范围”也不等于其中每个小节每年都出题。

\begin{center}
\begin{tabularx}{\textwidth}{@{}p{.38\textwidth}XX@{}}
\toprule 本书内容&数学一&数学三\\\midrule
第1—5章基础部分&按大纲读&按大纲读\\
第6章抽样及三大分布&按大纲读&按大纲读，注意典型构造\\
第6章经验分布函数&按本卷要求选读&大纲明确列入\\
第7.1节矩估计与MLE&按大纲读&按大纲读\\
第7.2节估计量评价&无偏、有效、一致等按要求读&不将全节列为共同必修\\
第7.5节区间估计&正态总体相应部分&不作本卷主线\\
第8章假设检验&检验概念与正态总体相应部分&不作本卷主线\\
第9章回归与ANOVA、第10章&提高选读&提高选读\\
\bottomrule
\end{tabularx}
\end{center}

充分完备、UMVU、信息界、Bayes、Delta 方法、测度论等按提高材料使用；不能因为它们位于基础章里，就与整个基础章一起列为统考核心。数学一的正态区间与检验也不是第7—8章所有推广题的统称。

\section*{三个已有题型的复习顺序}
\begin{enumerate}
\item \textbf{二维变换先画区域：}见第\ref{sec:33}—\ref{sec:34}节。先画联合密度支撑，明确变换不等式与零概率分母，按参数位置分段确定积分限；算出分布函数后再检查两端极限、单调性与总概率。
\item \textbf{抽样分布先列独立性：}见第\ref{sec:62}—\ref{sec:63}节。识别标准正态、独立卡方及相应自由度后再拼成 $t$ 或 $F$ 统计量。只看外形、没有独立性条件，不足以判断分布。
\item \textbf{MLE 先保留支持：}见第\ref{sec:71}节例\ref{ex:7-3}。均匀分布端点例题已经给出可行参数与样本最大值的关系，应与矩估计、正态及指数例题一起练习。无内点驻点不等于无边界最大值；点态密度版本和最大值是否取到也要核对。
\end{enumerate}

边界MLE的分类与非唯一情形见第\ref{r2:mle}节；边缘正态不推出联合正态的反例见例\ref{ex:3-15}。附录\ref{app:exam_core}提供基础题型入口，全文索引仍保留提高内容，便于后续课程使用。
范围依据为教育部教育考试院《2027年全国硕士研究生招生考试数学考试大纲》原书第20—25、45—49页。\href{https://xuanshu.hep.com.cn/front/book/findBookDetails?bookId=6aa1910be119ac97297a608f}{官方版本入口}。
